\section*{Architecture génératrice et conservation de l'information}
\addcontentsline{toc}{section}{Architecture génératrice et conservation de l'information}
\label{libnar:troncal}

L'holographie modulaire triadique, abrégée HMT, situe l'origine de la construction dans un système fini de positions, d'opérations et de transports. De son prolongement, on obtient des histoires compatibles ; des histoires, des lectures numériques et d'incidence ; de ces réalisations, les structures avec lesquelles opère la mécanique dimensionnelle. La question qui organise cet article concerne cette continuité constructive : ce qui demeure lorsqu'une unité change d'échelle, lorsqu'une opération modifie sa représentation et lorsqu'une partie de l'état est transférée à la mémoire.

Le contenu de l’unité se précise tout au long du parcours. Dans une
partition arithmétique, la somme normalisée de ses composantes est conservée.
Dans un espace muni d’un produit scalaire, les normes et les corrélations sont conservées.
Dans une représentation géométrique, les formes, les directions et les
incidences sont transportées. Chaque conservation possède un objet et une équation propres.
Leur articulation permet de suivre une même construction depuis la retenue
d’un chiffre jusqu’à la récupération d’un état distribué dans une archive
de profondeur arbitraire.

L’exposition suivante présente ce parcours dans son ordre de dépendance.
Les définitions et démonstrations complètes apparaissent ensuite, dans le
corps mathématique. Il importe ici que le lecteur puisse reconnaître le rôle
de chaque opération avant d’étudier sa formulation générale : comment se
produit une histoire, ce qui détermine une lecture et de quelle manière peut
être reconstruite l’information qu’une représentation visible a cessé
de montrer.

\subsection*{Des positions finies aux histoires opératorielles}
\label{libnar:app}

La construction nommée APP, pour \emph{All Possible Paths}, part de deux axes de neuf marques et de deux opérations sur leurs paires de positions : somme et produit. Les positions possèdent une identité qui est conservée lors de leur évaluation. La paire \((i,j)\) admet la lecture additive \(i+j\) et la lecture multiplicative \(ij\) ; un chemin enregistre en outre comment on est arrivé à cette paire, quelle opération a agi et quelle orientation a suivie le transport. La relation entre accumulation additive et composition multiplicative se trouve ainsi située sur un support commun.

L’étape modulaire conserve aussi bien le reste que le quotient. Pour une
évaluation positive \(n\), la marque nonadique \(r\), comprise entre
un et neuf, et le quotient \(q\) vérifient
\[
 n=r+9q.
\]
Cette égalité permet de récupérer l’évaluation entière à partir de sa lecture
modulaire enrichie. Deux positions peuvent partager la même marque et
conserver des quotients différents. Par exemple, \((5,5)\) et \((8,2)\)
produisent la somme dix ; leurs produits sont vingt-cinq et seize. Les deux
produits ont pour marque sept, avec des quotients deux et un, respectivement.
La différence opératorielle est localisée dans le registre qui accompagne
la marque. Les définitions de la \cref{sec:nucleo} développent cette
conservation conjointement avec les feuillets d’opération et les routes.

L’unité ternaire orientée, appelée TRIT, distingue trois régimes
au moyen des signes \(-1,0,+1\). Son action comprend le choix de phase,
l’orientation des prolongements et la retenue. La division ternaire
équilibrée écrit un entier sous la forme \(n=r+3c\), avec
\(r\in\{-1,0,+1\}\) ; le reste identifie l’état local et le
quotient conserve ce qui est transporté au niveau suivant. Les réalisations
géométriques ultérieures distinguent les régimes elliptique, parabolique et
hyperbolique. Dans la lecture temporelle du corpus, le retour avec mémoire,
le seuil et l’ouverture bidirectionnelle possèdent des actions spécifiques sur
l’état.

Le \emph{Topological Phase Kernel}, ou noyau topologique de phase, est abrégé TPK. Il compose la sélection tritique, le transport et la mise à jour de mémoire. Une transition modifie les positions, les feuillets ou la phase, et conserve les données nécessaires pour relier l'état précédent au suivant. Son résultat est un état enrichi : il contient des résidus, des quotients, des retenues, une orientation, une frontière et des préfixes de chemin. Le nombre qui sera lu ultérieurement provient de cette histoire opératorielle. Sa généalogie conserve les opérations qui permettent de le poursuivre et les relations qui permettent de le comparer à d'autres histoires.

Cette organisation donne un sens précis au retour. Le calendrier
à cent huit positions se répartit en douze fenêtres de neuf. Lorsqu’un
tour nonadique s’achève, la phase correspondante revient et le
registre de fenêtre avance. Le retour à une position visible et la
récupération de l’état enrichi sont donc deux événements
différenciés. La \cref{sec:extension} formalise le calendrier,
l’holonomie et l’accroissement de mémoire qui accompagne leurs retours.

\subsection*{La contraction de la valeur et la permanence de l’incidence}
\label{libnar:continuo}

Une histoire se prolonge au moyen de préfixes compatibles. Chaque préfixe
conserve les opérations effectuées et délimite ses continuations
admissibles. La construction conjointe du continu maintient dans ce même
système le transport gradué, les bilans de mémoire, les incidences
ternaires, les relations de frontière et les compositions
déterminantales. Ces opérations sont reliées par leurs actions sur
les mêmes histoires et par leurs compatibilités avec le raffinement
(\cref{iv:continuo-conjunto}). Leur corrélation fait partie de
la construction de l’objet.

La lecture archimédienne associe aux préfixes des cylindres de précision croissante. Leurs intervalles compatibles se contractent et déterminent une valeur lorsque leurs diamètres tendent vers zéro. La convention semi-ouverte distingue les représentants de frontière pendant le raffinement ; l'information de retenue enregistre les continuations qui atteignent une même extrémité. Dans cette lecture apparaît la réalisation numérique du continu engendré. La suite d'intervalles exprime dans quelle mesure la valeur a été déterminée ; l'état enrichi exprime comment elle a été déterminée.

La lecture incidentielle conserve une autre information : quels éléments forment
un bloc, comment une frontière est orientée, quels transports relient les
préfixes et quelles actions préservent ces relations. À travers les applications
développées dans la \cref{exc:seccion}, cette information se réalise dans
les matrices de Paley et de Hadamard, les plans de Witt et les systèmes de Steiner,
les codes de Golay, les constructions réticulaires et le réseau de Leech.
Le prolongement à l’algèbre d’opérateurs de vertex \(V^\natural\)
y situe l’action du groupe Monstre. Chaque étape conserve son porteur,
son application et ses conditions de compatibilité.

La relation entre nombre et symétrie devient ainsi concrète. La valeur
archimédienne identifie des histoires sous une équivalence numérique ;
l’incidence conserve des relations entre les histoires et leurs frontières.
Les deux lectures proviennent d’un même état et remplissent des fonctions
complémentaires. Une grandeur scalaire peut rester fixe tandis que
la représentation incidencielle subit une transformation. Suivre
les deux lectures permet d’étudier conjointement la stabilisation de la
valeur et l’évolution de la structure qui la réalise.

L’unité se conserve également lorsqu’elle se distribue entre les prolongements.
Les mesures des cylindres successeurs ont pour somme la mesure du cylindre prédécesseur.
Dans les espaces de fonctions associés, cette égalité devient
une isométrie de raffinement : une fonction du niveau précédent conserve
sa norme lorsqu’elle est représentée sur les nouvelles positions. Le même
principe relie partition, mesure et transport. Les démonstrations
de la construction conjointe spécifient la composition et le passage à la
limite, de sorte que chaque nouvelle profondeur reste liée aux
précédentes.

La réalisation d’une mesure ajoute une opération concrète à cette relation
entre état et lecture. Le système et le dispositif se couplent, produisent
un registre et établissent quelles distinctions deviennent accessibles. Dans une
réalisation isométrique du couplage, la description conjointe conserve
les corrélations ; le résultat enregistré détermine ensuite
l’actualisation correspondante. Évaluer un observable, conditionner
l’état à un résultat et préparer de nouveau le registre ont des fonctions
différenciées. Cette séquence permet de suivre tant l’information qui
parvient à la lecture que celle qui demeure dans l’état et dans le dispositif.
La mémoire intervient dans cette continuité opératoire entre une observation
et les suivantes (sections~\ref{nuclear:3} et~\ref{lib04:archivo}).

\subsection*{L’unité dimensionnelle et la retenue récupérable}
\label{libnar:unidad}

L’extrême et moyenne raison holographique est étudiée ici à travers les
opérations qui conservent l’unité au cours de son extension et de son
raffinement. La position ajoutée possède une identité propre. Si
\(E_9\) contient les neuf classes nonadiques, sa complétion s’écrit
\(E_9\sqcup\{\ast\}\). La classe résiduelle zéro appartient déjà à
\(E_9\) et a pour représentant positif neuf ; \(\ast\) désigne
l’état supplémentaire de frontière. La dixième possibilité conserve ainsi une
fonction distincte de la réduction modulaire. En plusieurs coordonnées, compter
combien ont atteint \(\ast\) détermine les strates de la complétion
et permet de suivre la contribution de chacune à l’unité.

La complétion cubique offre une première représentation exacte. En faisant passer de neuf à dix les positions de chaque coordonnée, l'unité cubique se décompose en intérieur, strates d'extension et sommet final :
\[
 10^3=9^3+3\cdot9^2+3\cdot9+1
     =729+243+27+1=729+271.
\]
Diviser par mille transforme le décompte en une partition de masse unité. Les termes conservent leur provenance dimensionnelle : une, deux ou trois coordonnées occupent la position ajoutée. Le développement \((9+1)^d\) conserve cette distinction pour chaque dimension entière positive. L'augmentation de dimension et le changement d'échelle sont reliés par des applications de coordonnées et des sommes de strates (\cref{sec:revision-emrh,lib03:completacion}).

La retenue introduit une transformation plus fine. Considérons la
partition \(73+27=100\). La multiplication par dix produit
\(730+270=1000\). Transférer une unité de la première composante à
la seconde donne
\[
 (73,27,1)\longmapsto(10\cdot73-1,10\cdot27+1)
                  =(729,271).
\]
La troisième donnée est la retenue. La somme finale reste mille et la
partition normalisée reste un. En même temps, le résidu décimal
de la seconde composante conserve le transfert : à partir de \(271\), on
retrouve le reste un, puis \((271-1)/10=27\), et enfin
\((729+1)/10=73\). Le résultat porte avec lui les données de
l’opération qui l’a produit.

Le théorème~\ref{lib03:teo-residuo} démontre cette inversion pour l’application
\[
 L_c(x,y)=(Bx-c,By+c),\qquad x+y=M,
\]
avec la base \(B\), la retenue canonique \(0\leq c<B\) et les conditions
de positivité déclarées. Après plusieurs étapes, le registre accumulé
\(C_n\) réunit les retenues comme un mot en base \(B\), et les
composantes satisfont
\[
 x_n=B^n x_0-C_n,\qquad y_n=B^n y_0+C_n.
\]
La somme conserve \(B^nM\). La base et la profondeur étant fixées, l'état final récupère le mot complet de retenues et la partition initiale (\cref{lib03:cor-palabra}). La profondeur fait partie de la donnée de reconstruction, car elle détermine également les positions initiales qui contiennent des zéros.

Les lectures normalisées convergent avec des bornes explicites pour la queue.
Trois niveaux de description sont ainsi reliés : une opération entière
réversible, une suite de partitions de l’unité et un système de
cylindres qui détermine une lecture limite. Le choix effectif des
continuations suit les règles du transport considéré. L’exemple
\((729,271)\) montre son premier transfert ; les preuves de
raffinement décrivent également la mémoire de tous les transferts
ultérieurs et les coïncidences de frontière entre développements
(\cref{lib03:teo-limite,lib03:teo-cilindros}).

\subsection*{Les régions numériques et le registre de couplage}
\label{libnar:registro}

Dans le corpus, les coordonnées \(\pi_{\rm HMT}\), \(\varphi_{\rm HMT}\) et \(e_{\rm HMT}\) proviennent de lecteurs régionaux sur le prolongement APP--TRIT--TPK. La région de clôture, la propagation et la relation stable entre modes possèdent leurs propres opérateurs. Par exemple, le calendrier des modes additif et multiplicatifs produit la matrice d'incidence \(F_{\rm av}=\bigl(\begin{smallmatrix}0&1\\1&1\end{smallmatrix}\bigr)\).
Sa direction positive stable détermine la coordonnée d’or dans la
réalisation correspondante. Les preuves de génération précisent, pour
chaque profondeur demandée, comment on obtient un préfixe certifié et
comment sa compatibilité avec les suivants est conservée
(\cref{sec:generacion}).

Le même état terminal conserve un registre orienté de ses douze
fenêtres. L’extraction des canaux signés, la charge totale et
l’inversion de Hadamard produisent le vecteur ordonné \(K\). Le terme
\emph{constante holographique de couplage arithmético-géométrique}
désigne ce registre avec ses lecteurs et avec l’origine de phase fixée.
Sa fonction consiste à relier un encodage arithmétique exact
aux réalisations d’incidence et de géométrie qui agissent sur le
même objet (\cref{lib01:definicion}).

Dans la représentation entière canonique, on obtient
\[
 \begin{aligned}
 K=\bigl(&234,543,140,729,659,824,\\
         &621,58,914,794,146,601\bigr),
 \qquad\sum_{j=1}^{12}K_j=6263.
 \end{aligned}
\]
La \cref{lib01:registro} montre l'extraction complète, les congruences de ses canaux et la divisibilité qui permet d'inverser la matrice de Hadamard. L'ordre de ses composantes conserve la position de fenêtre. Cette information intervient aussi bien dans la lecture des chiffres que dans les opérateurs cycliques et les incidences.

La somme totale remplit une fonction de reconstruction. Les différences
entre positions séparées de trois et quatre pas déterminent
conjointement les variations relatives du registre. Ajouter une même
quantité aux douze composantes conserve ces différences ; la somme fixe
précisément ce niveau commun. Dans \(K\), la charge \(Q(K)=6263\)
restitue la composante uniforme. Les différences, la charge et les compatibilités
entières permettent de retrouver le registre complet
(\cref{prop:k-transversal}). La relation globale entre les positions et
l’organisation relative de ses canaux sont conservées par des lectures
complémentaires.

En base mille, avec douze positions déclarées, le registre est codé par la fraction périodique
\[
 \kappa_{\rm ag}=
 \frac{234543140729659824621058914794146601}{10^{36}-1}.
\]
Chaque composante occupe trois chiffres ; c’est pourquoi la composante cinquante-
huit s’écrit \(058\) dans le développement. Multiplier la fraction
par son dénominateur et effectuer douze divisions euclidiennes récupère le
registre exact. La périodicité de cet encodage décrit une
lecture fixée du registre fini ; l’actualisation de l’histoire
enrichie conserve pour sa part sa propre dynamique de mémoire.

Le lien avec la coordonnée \(\alpha_A\) s'exprime dans le lecteur de clôture par
\[
 \alpha_A=
 \pi_{\rm HMT}+e_{\rm HMT}-\varphi_{\rm HMT}-4-\kappa_{\rm ag}.
\]
Les trois coordonnées régionales et le registre proviennent de l'état engendré et agissent dans l'ordre constructif décrit dans \cref{lib01:clausura}. La formule rend visible leur couplage arithmétique : l'extraction de \(K\) conserve une partie de la structure qui intervient dans cette clôture. Les autres réalisations de \(K\) conservent leurs lecteurs respectifs. Ainsi, une identité scalaire, une matrice d'incidence et un transport de phase peuvent être étudiés à partir de la même provenance, avec des actions clairement différenciées.

\subsection*{Un registre capable d’identifier une transformation}
\label{libnar:identificacion}

La fonction du registre peut être examinée au moyen d’une opération
concrète. Soit \(S\) le décalage cyclique de ses douze
composantes. Les douze colonnes
\(K,SK,\ldots,S^{11}K\) forment une matrice \(O_K\) inversible.
Le résultat signifie que l’orbite cyclique de ce registre atteint
toutes les directions de son espace à douze composantes. La preuve
calcule le spectre du Gram \(O_K^*O_K\) et donne des bornes positives
exactes pour ses valeurs propres (\cref{lib01:marco-ciclico}).

Supposons qu’un opérateur \(H\) respecte le décalage cyclique,
c’est-à-dire \(HS=SH\). Lorsqu’on observe la réponse vectorielle \(y=HK\),
ses décalages fournissent les réponses à toutes les colonnes
de \(O_K\). On récupère alors l’opérateur complet au moyen de
\[
 H=O_yO_K^{-1},\qquad
 O_y=(y,Sy,\ldots,S^{11}y).
\]
L’observation contient douze composantes, et l’hypothèse de
commutation identifie la classe des transformations récupérables.
Dans cette classe, une seule réponse vectorielle détermine tous
les coefficients du filtre circulaire. Par exemple, la combinaison
\(2I+3S-S^5\) est identifiée à partir de son action sur \(K\),
avec les bornes d’erreur du théorème~\ref{lib01:operadores}.

Ce fait offre une application opératoire de la structure du registre. Son ordre et son spectre permettent de l'utiliser comme état test pour identifier des transformations compatibles. Si l'état et ses réponses sont transportés au moyen d'une isométrie complète de mémoire, la matrice de Gram de l'orbite est conservée et l'identification reste récupérable par l'adjoint. La stabilité provient de valeurs propres calculées et d'opérateurs d'inversion explicites. Le même registre qui participe à la clôture arithmétique acquiert ainsi une fonction précise dans la reconnaissance des transports.

La normalisation possède également une information associée. Le facteur
unitaire de \(O_K\) conserve l’orientation et les angles ; son Gram
conserve les amplitudes nécessaires pour reconstruire la matrice
originale. La paire formée des deux facteurs retrouve \(O_K\)
(\cref{lib01:polar}). Cette séparation permet de choisir une
représentation normalisée tout en maintenant l’échelle
exacte du registre.

\subsection*{Deux transports et une conservation bilatérale}
\label{libnar:bilateral}

L’étape suivante relie le raffinement à une loi de
conservation des corrélations. Soient \(\mathcal B\) et
\(\mathcal C\) deux isométries qui transportent le même espace
d’états vers un espace de sortie commun. Pour \(a>0\), on forme les
deux opérateurs
\[
 Q_-=a\mathcal B-\mathcal C,\qquad
 Q_+=(a+1)\mathcal B+\mathcal C.
\]
Les deux combinaisons réunissent une contribution commune et une
contribution relative de signes différents. En additionnant leurs normes
avec les poids correspondants, les termes croisés s’annulent :
\[
 (a+1)\|Q_-v\|^2+a\|Q_+v\|^2
   =(2a+1)(a^2+a+1)\|v\|^2.
\]
Le théorème~\ref{lib04:teo-balance} démontre cette égalité pour toute dimension des espaces et pour les transports isométriques déclarés. La même annulation conserve les produits scalaires et admet sa formulation pour des formes préservées par les transports.

Dans le cas \(a=8\), le facteur \(a^2+a+1\) vaut soixante-treize.
La factorisation de l’isométrie complète distribue d’abord la
norme entre deux composantes de poids
\[
 \frac{72}{73}\qquad\text{et}\qquad\frac{1}{73}.
\]
Ensuite agit un mélange orthogonal qui produit les canaux
normalisés de \(Q_-\) et de \(Q_+\). Les poids ont ainsi une
origine opératorielle précise : \(72=8\cdot9\) et
\(73=8\cdot9+1\). Leur fonction se conserve lors du changement de dimension
ou des transports, dans les hypothèses de la loi bilatérale
(\cref{lib04:factorizacion}).

La récupération se voit directement dans les combinaisons de sortie. Si \(x=Q_-v\) et \(y=Q_+v\), alors
\[
 x+y=(2a+1)\mathcal Bv,\qquad
 ay-(a+1)x=(2a+1)\mathcal Cv.
\]
Les deux transports de l’état initial restent déterminés par
la paire de canaux. L’adjoint de chacune des isométries
retrouve ensuite l’état correspondant. Avec la normalisation
complète, l’opérateur \(W_a\) vérifie \(W_a^*W_a=I\) :
il conserve la norme et offre un inverse exact sur son image.

Une lecture particulière peut varier tandis que toute cette information est conservée. Lorsque \(\mathcal C=\mathcal B R\), avec \(\mathcal B\) et \(R\) unitaires, ce dernier comme transport relatif, la lecture diagonale d'un observable \(G\) dans les deux canaux produit
\[
 \mathcal T(G)=\frac{72G+R^*GR}{73}
 \qquad(a=8).
\]
Les observables compatibles avec \(R\) demeurent fixes ; les autres
enregistrent la différence de leurs deux actions. Cette équation distingue
la conservation de l’état complet et la réponse d’un lecteur
déterminé (\cref{lib04:teo-lector}).

\subsection*{La même retenue, désormais comme état transportable}
\label{libnar:ejemplo}

La relation précédente s'applique matériellement à la retenue. À chaque étiquette admissible \((x,y,c)\) est associé un vecteur d'une base orthonormale. La transformation entière \(L_c\) induit alors l'opérateur \(J\) qui envoie cette base sur les étiquettes de sortie \((Bx-c,By+c)\). Son inverse récupère la retenue au moyen du résidu. La bijection des étiquettes fait de \(J\) une transformation unitaire (\cref{lib05:prop-unitaria}). Les valeurs entières désignent les étiquettes ; les normes correspondent aux amplitudes des vecteurs qui les représentent.

Pour la base dix et la somme initiale cent, considérons les deux étiquettes
\(d_1=(73,27,1)\) et \(d_2=(73,27,2)\). Leurs images sont
\(e_1=(729,271)\) et \(e_2=(728,272)\). Fixons le transport
\(R\) qui échange ces deux étiquettes et laisse les autres fixes.
La composition utilise \(\mathcal B=J\) et \(\mathcal C=JR\).
Ce choix fournit un exemple explicite au sein du théorème
bilatéral, avec une règle de transport spécifiée avant son évaluation.

Sur l’état \(|d_1\rangle\), les canaux pour \(a=8\)
sont
\[
 Q_-|d_1\rangle=8|e_1\rangle-|e_2\rangle,
 \qquad
 Q_+|d_1\rangle=9|e_1\rangle+|e_2\rangle.
\]
Leurs normes au carré sont soixante-cinq et quatre-vingt-deux. Le bilan
exact est \(9\cdot65+8\cdot82=1241=17\cdot73\), tandis que
la somme des deux canaux retrouve \(17|e_1\rangle\).
Appliquer l’inverse des étiquettes retrouve la partition \((73,27)\)
et la retenue un.

La lecture fractionnaire de la première composante offre en outre une
prédiction exacte pour ce transport déclaré. En l’évaluant
diagonalement dans les deux canaux normalisés, on obtient
\[
 \frac{72\cdot729+728}{73\cdot1000}
   =\frac{729}{1000}-\frac{1}{73000}.
\]
Le lecteur diagonal de retenue donne \((72\cdot1+2)/73=74/73\).
L’observable transporté depuis l’entrée, en revanche, conserve
la valeur propre un de la retenue originale. La différence entre les deux
résultats permet d’identifier quelle opération de lecture a été
réalisée. L’information de l’état demeure récupérable et la
lecture diagonale quantifie comment le transport relatif intervient
(\cref{lib05:ejemplo}).

L’exemple réunit les deux conservations de l’article. La somme
arithmétique est conservée dans les étiquettes raffinées et la norme est
conservée dans leurs représentations unitaires. Le lien entre les deux
provient de la bijection et de son relèvement aux bases orthonormales.
Cette composition permet de calculer des lecteurs, d’inverser le transport
et de contrôler l’erreur tout en maintenant identifiés les objets sur
lesquels agit chaque loi.

\subsection*{Mémoire complète, symétries et profondeur arbitraire}
\label{libnar:memoria}

La conservation bilatérale se prolonge étape par étape. À chaque niveau, un canal poursuit l'état et l'autre enregistre un complément de mémoire. Si \(R_Nv\) est l'état qui atteint le niveau \(N\) et \(m_0,\ldots,m_{N-1}\) sont les compléments enregistrés, le bilan cumulé est
\[
 \|v\|^2=\|R_Nv\|^2+\sum_{n=0}^{N-1}\|m_n\|^2.
\]
L’égalité identifie exactement où demeure la norme
initiale. La récupération finie utilise l’état continué et
tous les compléments disponibles ; leurs corrélations sont conservées
par la même isométrie conjointe.

Pour la politique d'archivage du théorème~\ref{lib04:teo-archivo}, la norme du canal poursuivi possède une borne géométrique uniforme. Dans le cas \(a=8\), le carré de sa norme est borné à chaque étape par \(729/1241\). Lorsque la profondeur augmente, la composante poursuivie tend vers zéro et l'archive infinie conserve toute la norme. Si \(A_n\) et \(D_n\) désignent les blocs de continuation et de mémoire, respectivement, l'état se reconstruit par la série convergente
\[
 v=\sum_{n\geq0}R_n^*D_n^*m_n.
\]
L’opérateur d’inversion est de norme un. Une perturbation de norme
\(\varepsilon\) dans l’archive complète produit une erreur de
reconstruction de norme au plus égale à \(\varepsilon\). Les
troncatures finies possèdent leurs propres bornes et, lorsqu’on utilise
uniquement leur mémoire, l’inverse fini spécifié dans le même
développement.

L’incidence du registre \(K\), son orbite cyclique et ses
directions distinguées peuvent être transportées au moyen de ces
isométries. Les produits scalaires qui déterminent
les angles et les matrices de Gram sont conservés ; les projecteurs et les lecteurs sont
transportés avec leurs domaines ; l’action de l’adjoint récupère la
représentation initiale. Les puissances extérieures transportent
l’orientation et conservent les déterminants de Gram, qui mesurent les
carrés des volumes correspondants. La dimension du
porteur peut être étendue et l’information géométrique de
l’image demeure spécifiée par l’isométrie.

La relation avec la conservation de Noether est formulée à un niveau
supplémentaire bien déterminé. Pour l’action discrète auxiliaire
construite dans la \cref{lib01:carga-noether}, une collection de
générateurs qui entrelace les transports produit une charge
conservée de la forme \(p_n^*A_nq_n\). Interviennent ici une
action variationnelle, ses équations et la symétrie qui les préserve.
Le bilan de norme et la charge variationnelle s’articulent au moyen
des transports communs, chacun avec ses hypothèses et la
fonction qu’il remplit dans la démonstration.

Le parcours établit une relation opératoire entre unité, nombre, mémoire et géométrie. L'unité arithmétique admet des raffinements réversibles ; l'état enrichi conserve le registre des opérations ; les réalisations numérique et d'incidence extraient des propriétés corrélées ; le couplage bilatéral maintient l'information et les corrélations pendant leur redistribution. L'extrême et moyenne raison holographique acquiert ainsi un développement démonstratif dans lequel partition, extension, transport et récupération peuvent être suivis au moyen d'applications explicites. Les pages suivantes présentent ces applications, leurs preuves et leurs exemples complets.
